\section{The cobordism calculation}
\label{sec:cobordism}

The preceding construction gives candidate attaching classes and recovers
their unnormalized powers after inverting $t$.  To identify their
Hurewicz images, we now calculate these powers in complex cobordism.
Everything in this section takes place at one fixed odd prime $p$.
The relation we obtain will be applied, prime by prime, to the completed
towers in Section~\ref{sec:realization}.

Write $M_p=\MU_{(p)}$, and use cohomological grading for spaces, so that
$M_p^{-s}(\mathrm{pt})=(M_p)_s$.  We use the commutative multiplication
on the complex Thom spectrum and its complex orientation, supplied by
the coherent Thom construction
\cite[Chapter~IV, Section~2, Construction~2.5 and Lemma~2.2]{May1977}.
Localization at $p$ carries this structure to $M_p$: the smashing
localization $X\mapsto S_{(p)}\otimes X$ is idempotent, and its
equivalences are preserved by tensoring, so it is symmetric monoidal
\cite[Proposition~2.2.1.9]{LurieHA}. In particular,
complex bundles have multiplicative Thom classes, and tensor product of
line bundles defines a graded one-dimensional commutative formal group
law.  The coefficient calculation is
\[
 (M_p)_*=\Z_{(p)}[a_s:s\geq 1],\qquad |a_s|=2s;
\]
see \cite[Chapter~2, Section~2.4, Theorem~4 of the English translation]{Novikov1962}.
We also use Quillen's multiplicative projection
\[
 M_p\longrightarrow\BP
\]
and the standard $p$-typical orientation of $\BP$
\cite[Section~5, Theorem~4]{Quillen1969}.  A map of ring spectra in the
homotopy category is sufficient here.  We do not require this projection
to preserve power operations, and we do not place a commutative
ring-spectrum structure on $\BP$.

\subsection{Coefficients of the \texorpdfstring{$p$}{p}-series}

Let $Z$ be the Euler coordinate of the universal complex line, of
cohomological degree two.  Define homogeneous coefficients and ideals by
\begin{equation}
 \label{cob:p-series}
 [p](Z)=\sum_{u\geq 1}k_uZ^u,
 \qquad k_u\in(M_p)_{2(u-1)},
 \qquad x_h=k_{p^h},
 \qquad I_h=(x_0,\ldots,x_{h-1}).
\end{equation}
Here $I_0=0$, $x_0=p$, and $|x_h|=d_h=2(p^h-1)$.

\begin{lemma}
 \label{cob:regularity}
 The sequence $x_0,x_1,\ldots$ is regular in $(M_p)_*$.
 Under the multiplicative projection to $\BP_*$, the ideal $I_h$ has
 image $(p,v_1,\ldots,v_{h-1})$ for $h\geq 1$, and the image of $x_h$
 modulo this ideal is a scalar unit times $v_h$.  Moreover,
 \[
  k_u\in I_h\qquad(1\leq u<p^h).
 \]
\end{lemma}

\begin{proof}
 In the standard $p$-typical coordinate, the logarithm of the
 Brown--Peterson formal group is
 \[
  \log(Z)=\sum_{h\geq 0}\ell_hZ^{p^h},\qquad \ell_0=1,
  \qquad
  p\ell_h=\sum_{0\leq s<h}\ell_s v_{h-s}^{p^s}.
 \]
 This is the Hazewinkel recursion
 \cite[Section~3.1, Equations~(3.1.1)--(3.1.3)]{HazewinkelIII}.
 Over the rationals, the $p$-series is
 $\log^{-1}(p\log(Z))$.  Its indecomposable term at $Z^{p^h}$ is
 \[
  (p-p^{p^h})\ell_h
   =(1-p^{p^h-1})v_h
       \pmod{\text{lower generators}},\qquad h>0.
 \]
 By homogeneity no generator $v_j$ with $j>h$ occurs in this
 coefficient; every term other than the displayed multiple of $v_h$
 uses lower generators.  The coefficient itself is integral, so these
 remaining terms lie in $\Z_{(p)}[v_1,\ldots,v_{h-1}]$.
 The scalar $1-p^{p^h-1}$ is a $p$-local unit.  Also, modulo
 $(p,v_1,\ldots,v_{h-1})$, degree considerations exclude every
 coefficient of the $p$-series below $Z^{p^h}$.

 These leading-coefficient assertions are unchanged, up to a scalar
 unit, by a change of Euler coordinate.  Indeed, the projective-space
 formula identifies a second coordinate with an integral power series
 in the first, with invertible linear coefficient.  Conjugating a
 series whose terms below order $p^h$ vanish preserves that order and
 multiplies its leading coefficient by a power of this invertible
 linear coefficient.  For normalized complex orientations the change
 is strict, and the leading coefficient is unchanged.  Induction on
 $h$ therefore gives the asserted images of the ideals $I_h$ and of
 their next generators, also for the Euler coordinate induced from
 $M_p$.

 We next prove regularity in $(M_p)_*$ itself. For $h\geq1$, homogeneity in the
 polynomial coefficient ring gives
 \[
  x_h=c_h a_{p^h-1}
       +D_h(a_1,\ldots,a_{p^h-2}),\qquad c_h\in\Z_{(p)},
 \]
 where $D_h$ is decomposable.  Under the projection to $\BP_*$, the
 image of a generator of smaller degree cannot involve $v_h$.
 Consequently the coefficient of $v_h$ in the image of $x_h$ is the
 product of $c_h$ and the coefficient of $v_h$ in the image of
 $a_{p^h-1}$.  The preceding calculation says that this product is a
 unit modulo $p$.  Thus $c_h$ is a $p$-local unit.  Replacing
 $a_{p^h-1}$ successively by $x_h$ is a triangular change of polynomial
 generators.  Since $p$ is a nonzerodivisor, this proves regularity of
 $p,x_1,x_2,\ldots$.

 Finally, work modulo $I_h$, with $h>0$.  This coefficient ring has
 characteristic $p$.  Suppose that the $p$-series has a first nonzero
 term $cZ^s$.  Since it is a formal group endomorphism, comparison of
 total degree $s$ in its compatibility with formal addition gives
 \[
  c(X+Y)^s=cX^s+cY^s.
 \]
 Hence $\binom{s}{i}c=0$ for $0<i<s$.  If $s=p^rm$ with
 $p\nmid m$ and $m>1$, then
 $\binom{s}{p^r}\equiv m\pmod p$, which is a unit in any
 $\F_p$-algebra.  This is impossible for nonzero $c$.
 Thus $s$ must be a $p$-power.  All possible such coefficients below
 $p^h$ are among $x_0,\ldots,x_{h-1}$ and have already been killed.
 This proves $k_u\in I_h$ in the stated range.  The argument uses no
 assumption that the quotient ring is a domain.
\end{proof}

\subsection{Borel cohomology and Euler divisibility}

Retain $G=C_p\rtimes\F_p^\times$, its real permutation representation
$\rho$, and $V=\rho-1$.  For a subgroup $K\leq G$, set
$H_K^*=M_p^*(BK)$.  These are ordinary Borel cohomology groups; no
Euler inversion or categorical localization is being performed in this
section.

\begin{lemma}
 \label{cob:cohomology}
 For the Euler coordinate $x$ of a faithful complex character of $C_p$,
 restriction gives
 \begin{equation}
  \label{cob:borel-rings}
  H_{C_p}^*=(M_p)^*(\mathrm{pt})[[x]]/([p](x)),
  \qquad
  H_G^*\xrightarrow{\ \simeq\ }
     (H_{C_p}^*)^{\F_p^\times}.
 \end{equation}
 The projective-space formulas are
 \[
  M_p^*(BK\times\mathbb CP^r)=H_K^*[Z]/(Z^{r+1}),
  \qquad
  M_p^*(BK\times\mathbb CP^\infty)=H_K^*[[Z]],
 \]
 with graded power series in the infinite-dimensional case.
 Evaluation at a point of $BG$ detects units in $H_G^0$.
\end{lemma}

\begin{proof}
 The classes $1,Z,\ldots,Z^r$ form a free basis for
 $M_p^*(\mathbb CP^r)$.  To see this at the module level, the cellular
 spectral sequence collapses because $M_p$ has even coefficients.
 The class $Z$ is a generator in filtration two, and its powers
 generate the successive even filtrations by the ordinary cup-product
 calculation.  The resulting map
 \[
  \bigoplus_{j=0}^r\Sigma^{-2j}M_p
     \longrightarrow\map(\mathbb CP^r_+,M_p)
 \]
 of $M_p$-modules is an equivalence.  The class $Z^{r+1}$ vanishes
 because its filtration exceeds the dimension of $\mathbb CP^r$.
 Applying cochains from $BK$ to this finite splitting proves the
 product formula at finite $r$.  Passage to the inverse limit gives
 the power-series formula; the truncation maps are surjective, so
 there is no additional derived-limit term.

 The unit circle bundle of the $p$-th tensor power of the universal
 line over $\mathbb CP^\infty$ models $BC_p$.  Its sphere-bundle
 triangle, together with the Thom isomorphism, identifies the relevant
 cochain map with multiplication by $[p](x)$.  This power series is a
 nonzerodivisor in $(M_p)^*(\mathrm{pt})[[x]]$, since the coefficient
 ring is a polynomial domain.  The resulting exact sequence gives
 the first formula in Equation~\eqref{cob:borel-rings}.

 For the second formula, compute homotopy fixed points first by $C_p$
 and then by its complement $\F_p^\times$.  If a finite group $K$ has
 order invertible on a spectrum, homotopy fixed points calculate the
 invariants on homotopy groups without a boundedness hypothesis.
 Indeed, the map to the coinduction of the underlying spectrum and
 the sum map back have composite $|K|$.  Finite induction and
 coinduction agree, so dividing this composite by $|K|$ gives a
 retraction.  After taking homotopy fixed points, the coinduced term
 is the underlying spectrum, and the opposite composite is the
 averaging operator.  Apply this to $K=\F_p^\times$ and to the
 $C_p$-homotopy fixed points of the constant spectrum $M_p$.

 Finally, a degree-zero Borel class acts as an $M_p$-linear
 endomorphism of the constant rank-one $M_p$-module over $BG$.
 If its value at a point is a unit, this endomorphism is a fiberwise
 equivalence, hence an equivalence of local systems.  Its inverse is
 again such an endomorphism, proving the assertion about units.
\end{proof}

Define
\begin{equation}
 \label{cob:epsilon}
 \epsilon=e_{M_p}(V\otimes_{\mathbb R}\mathbb C)\in H_G^q,
 \qquad q=2(p-1).
\end{equation}
The restriction of this Euler class to $C_p$ is
\[
 \epsilon|_{C_p}=\prod_{\alpha=1}^{p-1}[\alpha](x)
                  =x^{p-1}\cdot\text{a unit}.
\]
Here $[\alpha](x)$ is the formal multiple for the character indexed by
$\alpha$, and its linear coefficient $\alpha$ is a $p$-local unit.

\begin{lemma}
 \label{cob:euler-divisibility}
 If a homogeneous class $y\in H_G^*$ restricts to an element of
 $(x^p)\subset H_{C_p}^*$, then $y$ is divisible by $\epsilon^2$ in
 $H_G^*$.
\end{lemma}

\begin{proof}
 For every $s>0$, the quotient $(x^s)/(x^{s+1})$ in $H_{C_p}^*$ is
 annihilated by $p$: multiply the relation $[p](x)=0$ by $x^{s-1}$.
 A multiplier $\alpha\in\F_p^\times$ acts on this quotient by
 $\alpha^s$, or its inverse if the opposite action convention is used.
 For
 \[
  p\leq s<2(p-1),
 \]
 this character is nontrivial.  An invariant element in $(x^s)$
 therefore maps to zero in $(x^s)/(x^{s+1})$, since some
 $\alpha^s-1$ is a unit modulo $p$.  Iterating over this finite range
 shows that $y|_{C_p}$ lies in $(x^{2(p-1)})$.

 Since $\epsilon|_{C_p}$ is $x^{p-1}$ times a unit, there is a
 homogeneous $z\in H_{C_p}^*$ with
 $y|_{C_p}=\epsilon^2z$.  Average $z$ over $\F_p^\times$.
 Both $y|_{C_p}$ and $\epsilon$ are invariant, so the averaged divisor
 still has product $y|_{C_p}$ with $\epsilon^2$.  By
 Lemma~\ref{cob:cohomology}, this invariant divisor comes from $H_G^*$.
 Notice that this argument neither asserts that $\epsilon$ is a
 nonzerodivisor nor cancels a factor of $\epsilon$.
\end{proof}

\subsection{Thom-corrected permutation powers}

We next define the ordinary primewise operation used in the
calculation. Its Thom construction belongs to the cobordism power
methods of tom Dieck and Quillen
\cite{tomDieck1968,Quillen1971}; the Euler factors will be retained
explicitly here. For a space $X$ and $f\in M_p^{2d}(X)$, represent $f$ by
a map
\[
 S^{-2d}\wedge X_+\longrightarrow M_p.
\]
Apply the permuted $p$-fold smash, pull back along the diagonal of
$X^p$, and use the commutative multiplication on $M_p$.  This gives a
map of spectra with naive $G$-action
\[
 S^{-2d\rho}\wedge X_+\longrightarrow\Tr M_p.
\]
The virtual real representation $2d(\rho-p)$ is complex.  Choose its
complex Thom trivialization
\[
 S^{2d(\rho-p)}\longrightarrow\Tr M_p.
\]
Tensoring the two maps and multiplying in $M_p$ replaces the source
representation by $-2pd$.  We obtain
\begin{equation}
 \label{cob:pi-definition}
 \Pi_p(f)\in M_p^{2pd}(BG\times X).
\end{equation}
The choices of Thom trivialization depend only on the integer $d$,
not on $X$ or $f$.  Such a Thom map induces an equivalence after
extension to $M_p$-modules because it does so on every fiber.  Negative
values of $d$ cause no difficulty: the corresponding trivializations
are obtained by dualizing those of actual complex bundles.  This
construction agrees with the even-degree Thom construction of
\cite[Section~3.2]{JohnsonNoel2010}, with allowance for the chosen
orientation units.

For a space $X$, let $\mathfrak t_G(X)$ denote the additive subgroup of
$M_p^*(BG\times X)$ generated by transfers from the subgroups of $G$
whose orders are prime to $p$.  It is a graded ideal by the projection
formula.  Write $\mathfrak t_G=\mathfrak t_G(\mathrm{pt})$.

\begin{lemma}
 \label{cob:power-rules}
 The operation $\Pi_p$ is natural in $X$.  For homogeneous even-degree
 classes $f$ and $g$, its product rule has the form
 \[
  \Pi_p(fg)=\alpha\,\Pi_p(f)\Pi_p(g),
  \qquad \alpha\in(H_G^0)^\times,
 \]
 where $\alpha$ depends on their degrees and the Thom choices and is
 pulled back from $BG$.  On any finite sum of classes in the same
 even degree, $\Pi_p$ is additive modulo $\mathfrak t_G(X)$.
 For $X=\mathbb CP^r$, extracting a coefficient of $Z^s$ sends
 $\mathfrak t_G(X)$ into $\mathfrak t_G$.
\end{lemma}

\begin{proof}
 Naturality follows from the diagonal construction.  Before the Thom
 correction, the product rule follows from the symmetric monoidal
 permuted smash and the commutative multiplication of $M_p$.
 Compare the product of the two chosen Thom trivializations with the
 chosen trivialization for the sum of the degrees.  Their ratio is
 an automorphism of a constant suspended rank-one $M_p$-module over
 $BG$, hence a unit in $H_G^0$.  This comparison is made before
 introducing $X$, so the factor has no dependence on its
 cohomology classes.

 To examine a finite sum of classes in one degree, factor the sum
 map through a finite biproduct and distribute the $p$-fold smash
 over that biproduct.  The summands are indexed by label tuples on
 the $p$ letters.  A constant label tuple gives exactly the power
 of the corresponding summand.  The stabilizer of a nonconstant
 tuple has order prime to $p$: otherwise it contains the unique
 subgroup $C_p$ of order $p$, whose transitive action on the letters
 would force the labels to be constant.  The other orbits therefore
 factor through inductions from prime-to-$p$ subgroups.

 A composite between trivially acting endpoints that factors through
 an induction from a subgroup is the corresponding additive transfer.
 This follows by the induction--restriction adjunction and the
 coinduction--restriction adjunction, using equality of finite
 induction and coinduction; their unit and counit give the transfer
 of the composite at the subgroup.  Tensoring with a Thom
 trivialization and multiplying do not change the property of
 factoring through such an induction.  This proves additivity
 modulo $\mathfrak t_G(X)$ and also explains directly its stability
 under multiplication by classes from the larger group.

 For the last assertion, use the finite $M_p$-module splitting of
 the cochains of $\mathbb CP^r$ from Lemma~\ref{cob:cohomology}.
 Coefficient extraction is the projection onto one of its summands.
 Taking cochains from a space with trivial group action commutes
 with finite induction, since that induction is a finite sum.
 Consequently these projections commute with transfers and take
 a transfer error to a transfer error on the coefficient spectrum.
\end{proof}

Let $P$ be the universal line bundle on $\mathbb CP^\infty$, and put
$h(Z)=\Pi_p(Z)$.  The next formula both identifies its first
coefficient and controls the remaining coefficients by the
Euler filtration.

\begin{lemma}
 \label{cob:euler-power}
 There is a unit
 $U(Z)\in M_p^0(BG\times\mathbb CP^\infty)^\times$ such that
 \begin{equation}
  \label{cob:euler-product}
  h(Z)=U(Z)\,e_{M_p}\bigl(P\otimes
                      (\rho\otimes_{\mathbb R}\mathbb C)\bigr).
 \end{equation}
 In particular, $h(Z)$ has no constant term and
 \begin{equation}
  \label{cob:linear-coefficient}
  [Z]h(Z)=U(0)\epsilon,
  \qquad
  h(Z)|_{C_p}=U(Z)|_{C_p}
          \prod_{\alpha=0}^{p-1}\bigl(Z+_{\MU}[\alpha](x)\bigr).
 \end{equation}
\end{lemma}

\begin{proof}
 Represent $Z$ by the zero section of $P$ followed by its Thom map.
 Fiberwise, the normalized Thom map is a map from the local sphere
 $S^{P-\mathbb C}$ to $M_p$ that becomes an equivalence after
 extending to $M_p$-modules.  Thom spectra can be formed as colimits
 of these sphere local systems.  Thus the external permuted smash
 of the Thom maps, which commutes with their colimits, is the
 corresponding fiberwise construction on the external sum of the
 bundles.  On pulling back to the diagonal and making the Thom
 correction in Equation~\eqref{cob:pi-definition}, the zero section
 becomes that of
 $P\otimes(\rho\otimes_{\mathbb R}\mathbb C)$.

 The subsequent map of its sphere local system, normalized by
 $\mathbb C^p$, is an $M_p$-linear equivalence after extension: on
 each fiber it is the product of the chosen units.  Its ratio to
 the ordinary Thom equivalence of this bundle is therefore a unit
 over $BG\times\mathbb CP^\infty$.  This is the class $U(Z)$ in
 Equation~\eqref{cob:euler-product}.

 Restricted to $C_p$, the complex permutation representation splits
 into the characters indexed by $\alpha\in\F_p$.  Multiplicativity
 of Euler classes and the formal group law give the product in
 Equation~\eqref{cob:linear-coefficient}.  The trivial summand gives
 the factor $Z$, so the constant term vanishes; evaluating every
 other factor at $Z=0$ gives the linear coefficient $U(0)\epsilon$.
 With compatible canonical Thom choices this is the usual exact
 Euler-product formula
 \cite[Equations~(4.11) and~(4.14)]{JohnsonNoel2010}; the unit-factor
 version suffices here.
\end{proof}

\subsection{Extraction of the next coefficient}

Fix an integer $a\geq0$. We now compare the power of the $p$-series
with the $p$-series substituted into $h$. The relation below has leading terms
$\Pi_p(x_a)$ and $\epsilon x_{a+1}$, each with a unit coefficient.
It separates four kinds of error: multiples of $\epsilon^2$,
products $x_jx_k$ with $j,k\leq a$, earlier powers $\Pi_p(x_j)$
with $j<a$, and additive transfers. In the Hurewicz induction of
Section~\ref{sec:realization}, the first two kinds will become
divisible by $t^2$, and the last two will vanish. This leaves the
comparison of leading terms from which one factor of $t$ can be
canceled before reduction modulo $t$.

\begin{proposition}
 \label{cob:coefficient-relation}
 For each integer $a\geq 0$ there are homogeneous coefficients in
 $H_G^*$ and units $C,C'\in(H_G^0)^\times$ such that
 \begin{equation}
  \label{eq:cobordism}
  \begin{split}
   C\Pi_p(x_a)={}&C'\epsilon x_{a+1}+\epsilon^2D_*
      +\sum_{0\leq j,k\leq a}D_{jk}x_jx_k\\
      &+\sum_{0\leq j<a}H_j\Pi_p(x_j)+T_*,
      \qquad T_*\in\mathfrak t_G.
  \end{split}
 \end{equation}
 All terms have cohomological degree $-p d_a$; coefficients in the
 two sums and $D_*$ are in the complementary degrees.
\end{proposition}

\begin{proof}
 Set $r_*=p^{a+1}$, and work on $\mathbb CP^{r_*}$.  Tensoring the
 universal line $p$ times gives a classifying map to
 $\mathbb CP^\infty$ that pulls $Z$ back to $[p](Z)$.  Naturality
 of $\Pi_p$ therefore gives $h([p](Z))$ for the power of this
 Euler class.  All substitutions can be made at finite truncation:
 by cellular approximation, the classifying map from
 $\mathbb CP^{r_*}$ factors up to homotopy through its
 $2r_*$-dimensional skeleton, and $Z^{r_*+1}=0$ on this space.

 Alternatively, expand $[p](Z)$ as in Equation~\eqref{cob:p-series}.
 Each summand $k_uZ^u$ has cohomological degree two, so the
 additivity assertion of Lemma~\ref{cob:power-rules} applies.
 Repeated use of its product rule gives
 \begin{equation}
  \label{cob:series-comparison}
  h([p](Z))\equiv
     \sum_{1\leq u\leq r_*}
          \alpha_u\Pi_p(k_u)h(Z)^u
       \pmod{Z^{r_*+1},\ \mathfrak t_G(\mathbb CP^{r_*})},
 \end{equation}
 where every $\alpha_u$ is a unit in $H_G^0$ pulled back from
 $BG$.  These factors may depend on $u$ and on the Thom choices;
 no compatibility between them is required below.

 Extract the coefficient of $Z^{r_*}$.  Since both sides of
 Equation~\eqref{cob:series-comparison} have degree $2p$, this
 coefficient has degree
 \[
  2p-2r_*=-p d_a.
 \]
 The transfer error remains an element of $\mathfrak t_G$ by
 Lemma~\ref{cob:power-rules}.

 First consider the term $u=p^a$ on the right.  Put
 $A_u=[Z^{r_*}]h(Z)^u$.  Its degree at $u=p^a$ is zero.
 At a point of $BG$, the permutation bundle is trivial and
 Equation~\eqref{cob:euler-product} becomes
 $h(Z)=U(Z)Z^p$.  Consequently
 \[
  A_{p^a}|_{\mathrm{pt}}=U(0)|_{\mathrm{pt}}^{\,p^a},
 \]
 a scalar unit.  Lemma~\ref{cob:cohomology} shows that
 $A_{p^a}$ is a unit of $H_G^0$.  Multiplication by
 $\alpha_{p^a}$ gives the coefficient $C$ of $\Pi_p(x_a)$.

 If $u>p^a$, restrict to $C_p$.  Each of the $pu$ formal-addition
 factors in the product expression for $h(Z)^u$ lies in the
 ideal $(Z,x)$.  The additional factor $U(Z)^u$ is an ordinary
 power series in these variables, and hence cannot decrease this
 ideal order.  Therefore
 \[
  A_u|_{C_p}\in(x^{pu-r_*}),\qquad pu-r_*\geq p.
 \]
 This assertion can be checked using any power-series lifts before
 quotienting by $[p](x)$; all exponents of $x$ and $Z$ are
 nonnegative.  Since $A_u$ comes from $BG$,
 Lemma~\ref{cob:euler-divisibility} gives
 $A_u\in\epsilon^2H_G^*$.  After multiplication by
 $\alpha_u\Pi_p(k_u)$ and summation, all such terms contribute
 to $\epsilon^2D_*$.

 For $1\leq u<p^a$, Lemma~\ref{cob:regularity} gives a finite
 homogeneous expression
 \[
  k_u=\sum_{0\leq j<a}c_{u,j}x_j
       \quad\text{in }(M_p)_*.
 \]
 All summands have the same degree as $k_u$.  Applying the sum
 and product rules of Lemma~\ref{cob:power-rules} expresses
 $\Pi_p(k_u)$, modulo transfers, as a sum of multiples of
 $\Pi_p(x_j)$ for $j<a$.  Multiplying by $\alpha_uA_u$ and
 collecting terms yields the last ideal sum in
 Equation~\eqref{eq:cobordism}.  When $a=0$ this range is empty.

 It remains to describe the coefficient on the left of
 Equation~\eqref{cob:series-comparison}.  Write
 $h(Z)=\sum_{s\geq 1}h_sZ^s$.  Its linear term contributes
 \[
  h_1[Z^{r_*}][p](Z)
      =U(0)\epsilon x_{a+1}.
 \]
 This is the displayed main term, with a unit coefficient.
 Every contribution from $h_s[p](Z)^s$ for $s\geq 2$ is a
 product of at least two coefficients $k_u$ with
 $1\leq u<r_*$.  Each is in $I_{a+1}$ by
 Lemma~\ref{cob:regularity}.  Thus these contributions belong to
 $I_{a+1}^2H_G^*$ and can be written as the displayed sum of
 $x_jx_k$ for $j,k\leq a$.

 Move the other right-hand terms to the left or absorb their signs
 into their coefficients.  The unit multiplying the main term
 remains a unit, giving Equation~\eqref{eq:cobordism}.  Every sum
 used above is finite at this fixed prime.  After collection there
 are only the displayed finitely many ideal factors, although
 their coefficients and the expression for $T_*$ may vary with
 $p$ without a uniform bound on the number of intermediate terms.
\end{proof}
